\section{Introduction}
An agent receives failing test output on a buggy draft, inspects the code, and prepares a fix. In Rich3518/R04 and FastAPI13786/R04, it formulates an exact, uniquely matching patch (\texttt{highlight=self.highlight} and \texttt{self.model.in\_.value}); in FastAPI14986/R03 and FastAPI15800/R04, it is still localizing the second file of a two-part change when its eighth call finishes. In three feedback traces (and two controls), the agent hits \texttt{BudgetExceeded} at event 8 and emits \texttt{edit\_file} at event 9---rejected because reading and editing share one pool, with 30--79 s left. Meanwhile, in Rich3067/R03, one branch lands its regex fix and passes, while its twin burns 17 \texttt{edit\_file} calls after delimiter corruption fuses \texttt{old\_string:} into \texttt{new\_string}.

Under realistic time and tool caps, upstream diagnostic gains are absorbed by downstream bottlenecks before a patch lands. We study this gap in Gemma 4 31B QAT across 3,517 runs and realizations on 95 tasks from 15 repositories: whole-workflow repair (736 runs; Study A), shared-draft continuations pairing ordinary (C) and advance-result (F) branches on 54 tasks (216 realizations, 648 endpoints; Study B), and three 855-run prompt-mitigation ablations (2,565 runs). We ask: \textbf{RQ1}, where does check construction fail between launch syntax, target reach, and grounding? \textbf{RQ2}, when advance results save check runs, how do agents reallocate freed budget, and what blocks diagnosed fixes from landing?
